\begin{statement}[C.8.13; Proved: the cusp exclusion]
If $\Im\tau\ge7/5$, then $|F|<17/20$.
\end{statement}
\begin{proof}
Here $|q|<1/6600$. On the auxiliary circle $|q|=1/2000$,
the Lambert series and eta product give
\[
 |E_6-1|\le
 504r\frac{1+26r+66r^2+26r^3+r^4}{(1-r)^7},\qquad
 |\Delta_0|\le r\exp\!\frac{24r}{(1-r)^2},
 \quad r=\frac1{2000}.
\]
The right sides imply $|\Delta_0/E_6^{3/2}|<1/1250$.
In particular $E_6$ has no zero on this disc and the specified
branch is holomorphic there. By Cauchy's inequality, with
$s=10/33$,
\[
 \sum_{n\ge1}\frac{|a_nq^n|}{n}
 \le\frac{-\log(1-s)}{1250},\qquad
 \sum_{n\ge1}\frac{|a_nq^n|}{n^2}
 \le\frac1{1250}\left(s+\frac{s^2}{4(1-s)}\right).
\]
Equation \eqref{eq:cc-EZ} therefore yields
\[
 |\mu-(1/2-i\beta)|<\frac1{125},\qquad
 |\mu'|<\frac{13}{250},\qquad \beta<\frac{365}{1000}.
\]
For $r=1/6600$, the corresponding weight-four and weight-six
Lambert bounds are
\[
 240r\frac{1+4r+r^2}{(1-r)^5}<\frac{37}{1000},\qquad
 504r\frac{1+26r+66r^2+26r^3+r^4}{(1-r)^7}
       <\frac{77}{1000}.
\]
Thus $|E_4|>963/1000$ and $|E_6|^{1/2}<1038/1000$.
Since $6\pi|p/g_2|=\sqrt6|E_6|^{1/2}/|E_4|$,
\begin{equation}\label{eq:cc-cuspbound}
 |F|<\sqrt6\frac{1038}{963}
 \left\{\frac{365/1000+1/125}{7/5}+\frac{13}{250}\right\}
 <\frac{17}{20}.
\end{equation}
All scalar comparisons here follow from geometric series and
the exponential Taylor bounds \eqref{eq:ca-exp}; for the logarithm
use $-\log(1-s)=\sum s^n/n$ with its geometric remainder.
For reference, the upper quantities before the last rounding are
less than $789/10^6$, $755/10^5$, $51/1000$, and $841/1000$
for $|\Delta_0/E_6^{3/2}|$, the error in $\mu$, the bound on
$\mu'$, and the final bound on $|F|$, respectively.
\end{proof}
